% Supplementary Table 6: wet-lab feedback, branch adjudication and locked recommendation.
\clearpage
\begingroup
\scriptsize
\setlength{\tabcolsep}{2pt}
\renewcommand{\arraystretch}{1.14}
\begin{longtable}{
  >{\raggedright\arraybackslash}p{0.13\linewidth}
  >{\raggedright\arraybackslash}p{0.15\linewidth}
  >{\raggedright\arraybackslash}p{0.14\linewidth}
  >{\raggedright\arraybackslash}p{0.19\linewidth}
  >{\raggedright\arraybackslash}p{0.17\linewidth}
  >{\raggedright\arraybackslash}p{0.13\linewidth}
}
\caption{Wet-lab feedback, branch adjudication and locked recommendation.}
\label{tab:feedback_locking}\\
\toprule
\textbf{Step} & \textbf{Input evidence} & \textbf{Question} & \textbf{Adjudication and caution} & \textbf{Decision or action} & \textbf{Lock status} \\
\midrule
\endfirsthead

\multicolumn{6}{l}{\tablename~\thetable\quad continued} \\
\toprule
\textbf{Step} & \textbf{Input evidence} & \textbf{Question} & \textbf{Adjudication and caution} & \textbf{Decision or action} & \textbf{Lock status} \\
\midrule
\endhead

\midrule
\multicolumn{6}{r}{Continued on next page} \\
\endfoot

\bottomrule
\endlastfoot

First-round feedback packet
& Returned IgG-binding reductions from matched donor, temperature, time-anchor and process-control branches.
& Which branches deserved refinement rather than broad expansion?
& Results were interpreted against untreated baseline and available process controls, not as a single leaderboard. Ultrasound was not isolated because no matched no-ultrasound result set was available.
& Retained donor ranking, temperature trade-off and pentose-entry questions; deferred broad wet-mode, ratio, water-activity and pretreatment expansion.
& Open adjudication. \\

Branch review
& Exact 4~h anchors showed the best tested hexose branch remained below the high-reduction target region; short-window pentose signals were stronger.
& Should exact-tested conditions be preferred over inferred pentose advancement?
& Exact support was treated as necessary for anchoring, but not sufficient for final promotion when the exact-tested branch had a low ceiling. Ribose was considered an inference-supported candidate, not a proven optimum.
& Advanced the pentose route for focused recommendation while recording uncertainty from the absence of exact ribose 60~$^\circ$C, 4~h data before validation.
& Candidate not yet locked. \\

Time-gain calibration
& Glucose 55~$^\circ$C, +US time course: 70.99\%, 78.75\%, 82.50\%, 84.33\%, 87.91\% and 90.81\% at 3, 6, 9, 12, 18 and 24~h.
& Did extra reaction time justify continued extension of the hexose route?
& Time increased reduction, but the trajectory was non-linear and implied rising process burden. The glucose curve was used only to calibrate diminishing marginal gains and was not treated as a ribose prediction curve.
& Avoided indefinite extension of the glucose branch; used the curve as a calibration source for evaluating shorter high-reactivity donor routes.
& Calibration evidence only. \\

Evidence hierarchy for final donor
& Exact anchors, ribose 3~h support, sequence-level priors from BepiPred-3.0 and NetGlycate, process burden and human feasibility constraints.
& Which donor had the strongest evidence-weighted path toward high IgG-binding reduction?
& Ribose had the strongest short-window pentose signal and a plausible chemistry-supported route beyond the exact-tested hexose range. Mechanistic reasoning was supportive, not site-specific proof; AGE and overprocessing risk were treated as secondary feasibility filters.
& Selected ribose as the final donor candidate under the fixed dry-state, ultrasound-assisted framework.
& Pre-lock recommendation. \\

Pre-validation lock
& Final agent adjudication from the locked wet-lab feedback record and Fig.~3c source data.
& What condition should be tested without post hoc alteration?
& The recommendation was recorded before validation. The estimate was heuristic and evidence-weighted, not a calibrated statistical interval.
& Locked ribose, 60~$^\circ$C, 4~h, ultrasound ON (400~W, 12~min), dry-state glycation, $a_w$ 0.79, protein:sugar 1:2 (w/w), 20~mM phosphate buffer at pH~7.3. Locked estimate: 86\% reduction, plausible range 82--93\%.
& Locked before validation. \\

Validation and confirmation
& Focused ribose 60~$^\circ$C, 4~h validation observed 85.6\% $\pm$ 1.7\% IgG-binding reduction; later donor-by-time panel was confirmatory.
& Did validation support the locked recommendation?
& The observed value was compared with the locked range only after recommendation lock. Agreement supported prioritization logic, not calibrated predictive accuracy. Later confirmatory data did not retroactively alter the recommendation.
& Treated the ribose condition as experimentally supported within the tested dry-state system and preserved the pre-validation decision record unchanged.
& Validation after lock. \\

\end{longtable}
\begin{flushleft}
\scriptsize IgG-binding reductions are expressed relative to the untreated $\beta$-LG baseline. ``Lock status'' records whether a row was part of open adjudication, calibration, pre-validation recommendation or post-lock validation. The table summarizes decision records and does not replace the temporarily held run-level result tables.
\end{flushleft}
\endgroup
